\documentclass[11pt,a4paper]{article}

\usepackage[margin=2.5cm]{geometry}
\usepackage{amsmath,amssymb}
\usepackage{booktabs}
\usepackage{longtable}
\usepackage{array}
\usepackage{enumitem}
\usepackage{hyperref}
\usepackage{xcolor}
\usepackage{listings}

\lstset{basicstyle=\ttfamily\small, breaklines=true}

\newcolumntype{L}[1]{>{\raggedright\arraybackslash}p{#1}}

\title{\textbf{IBM Module --- Complete Reference}\\[0.5ex]
  \large CERISSE \texttt{src/ibm}, BVH Backend}
\author{CERISSE Project}
\date{\today}

\begin{document}
\maketitle
\tableofcontents
\newpage

% ============================================================
\section{Code Layout}
% ============================================================

The IBM module is organized as a layered custom BVH stack integrated with AMReX.
The active geometry backend is the custom implementation in \texttt{bvh\_defs.h},
\texttt{bvh.h}, and \texttt{eib\_bvh.h}. Legacy CGAL material remains as reference
code in \texttt{eib\_cgal.h}, but the production IBM path is the BVH path.

\renewcommand{\arraystretch}{1.2}
\begin{longtable}{@{}L{3.5cm}L{10.5cm}@{}}
\toprule
\textbf{File} & \textbf{Role} \\
\midrule
\endfirsthead
\toprule
\textbf{File} & \textbf{Role} \\
\midrule
\endhead
\texttt{bvh\_defs.h}       & Core data definitions: \texttt{Point}, \texttt{Vec}, \texttt{BoundedSide}, \texttt{AABB}, \texttt{ClosestPointResult}, \texttt{BVHNode}, \texttt{TriMesh}/\texttt{Polygon2D}, \texttt{LocalFrame}, \texttt{SurfElem}. \\
\texttt{bvh.h}              & BVH algorithms: Morton encoding, primitive closest-point utilities, BVH build, BVH closest-point queries on CPU/GPU. \\
\texttt{eib\_bvh.h}         & IBM/BVH bridge: \texttt{InsideTester} + \texttt{InsideTesterView}, geometry readers (STL/OFF/polygon), cache builders, geometry consistency checks. \\
\texttt{eib\_data.h}        & IBM solver SoA data containers and constants: \texttt{gpData\_t}, \texttt{surfImp\_t}, \texttt{surfPhys\_t}, \texttt{FaceCSR}, interpolation thresholds and factors. \\
\texttt{eib.h}              & Main solver class template \texttt{eib\_t<wallmodel,param,cls\_t>} and runtime kernels: markers, ghost-point setup, GP reconstruction, surface reconstruction/output. \\
\texttt{eib\_io.h}          & Geometry-loading stage called by \texttt{eib\_t::init}: reads \texttt{ib.filename}, builds BVH, initializes in/out testers, builds flattened surface caches. \\
\texttt{eib\_interp.h}      & Interpolation/image-point search helpers used by GP and SURF reconstruction. \\
\texttt{IBMultiFab.h}       & IBM-specific FAB/MultiFab containers combining marker fields and per-FAB ghost-point SoA payload. \\
\texttt{ib\_walltypes.h}    & Wall model templates used by \texttt{compute\_surfIB} dispatch. \\
\texttt{eib\_geometry\_bvh.h} & Compatibility forwarding header $\to$ \texttt{eib\_bvh.h}. \\
\texttt{eib\_bvh\_geom.h}   & Compatibility forwarding header $\to$ \texttt{eib\_bvh.h}. \\
\texttt{bvh\_spatial\_index.h} & Compatibility forwarding header $\to$ \texttt{bvh.h}. \\
\bottomrule
\end{longtable}

% ============================================================
\section{CGAL to BVH Type Mapping}
% ============================================================

One-to-one type mappings between the original CGAL-based geometry stack and the custom
BVH-based implementation, focusing on equivalent functionality.
Source placement: \texttt{bvh\_defs.h} (data types), \texttt{bvh.h} (algorithms),
\texttt{eib\_bvh.h} (IBM/BVH bridge and geometry I/O), \texttt{eib\_data.h}
(solver SoA containers), \texttt{eib\_io.h} (geometry-loading flow).

\renewcommand{\arraystretch}{1.3}
\begin{longtable}{p{0.18\textwidth} p{0.24\textwidth} p{0.24\textwidth} p{0.26\textwidth}}
\toprule
\textbf{Functionality} & \textbf{CGAL Type} & \textbf{Custom BVH Type} & \textbf{Notes} \\
\midrule
\endfirsthead
\toprule
\textbf{Functionality} & \textbf{CGAL Type} & \textbf{Custom BVH Type} & \textbf{Notes} \\
\midrule
\endhead

Point representation
  & \texttt{Kernel::Point\_2 / Point\_3}
  & \texttt{GpuArray<Real, AMREX\_SPACEDIM>}
  & GPU-friendly fixed-size storage; same semantic role. \\ \midrule

Vector representation
  & \texttt{Kernel::Vector\_2 / Vector\_3}
  & \texttt{Vec = GpuArray<Real, AMREX\_SPACEDIM>}
  & Same geometric meaning (direction/difference vectors). \\ \midrule

Inside/outside enum
  & \texttt{CGAL::Bounded\_side}
  & \texttt{enum class BoundedSide}
  & Same three states: inside, boundary, outside. \\ \midrule

Bounding box type
  & \texttt{CGAL::Bbox\_2 / Bbox\_3}
  & \texttt{AABB} (alias \texttt{Bbox = AABB})
  & Equivalent role for pruning/containment queries. \\ \midrule

2D geometry container
  & \texttt{Polygon2D} (CGAL wrapper)
  & \texttt{Polygon2D} (custom)
  & Both represent polygon geometry in 2D. \\ \midrule

3D geometry container
  & \texttt{CGAL::Polyhedron\_3}
  & \texttt{TriMesh}
  & Triangle-only explicit indexing replaces halfedge structure. \\ \midrule

Spatial acceleration structure
  & \texttt{CGAL::AABB\_tree<Traits>}
  & \texttt{BVH}
  & Both accelerate nearest-point and ray-intersection traversal. \\ \midrule

Acceleration internals
  & \texttt{Primitive + Traits + Tree internals}
  & \texttt{BVHNode + BVH::nodes}
  & Custom stack exposes flat node arrays; CGAL hides internals. \\ \midrule

Closest-point query result
  & \texttt{Point\_and\_primitive\_id}
  & \texttt{ClosestPointResult}
  & Both return closest point + primitive id; custom type also stores distance. \\ \midrule

Primitive identifier
  & \texttt{Tree::Primitive\_id}
  & \texttt{int prim\_id}
  & Same purpose: identifies edge/face primitive hit by query. \\ \midrule

Element descriptor
  & \texttt{elm\_descriptor} (iterator / \texttt{face\_descriptor})
  & Direct integer indexing
  & Handle/descriptor abstractions replaced with linear ids. \\ \midrule

Inside tester functor
  & \texttt{Side\_of\_triangle\_mesh} (or 2D wrapper)
  & \texttt{inside\_t = InsideTester}
  & Unified callable interface retained: \texttt{operator()(Point)}.
    In 3D uses BVH-backed queries; in 2D uses BVH-accelerated edge culling. \\ \midrule

Primitive-id map
  & \texttt{std::map<PrimitiveID, int>}
  & Not required
  & Removed; direct integer indexing makes it unnecessary. \\ \bottomrule

\end{longtable}

% ============================================================
\section{Initialization Flow}
% ============================================================

\subsection{Entry Point and Call Chain}

At startup the first IBM call is \texttt{IBM::ib.init(\&amr)}.
The effective sequence is:

\begin{enumerate}[nosep]
  \item \texttt{eib\_t::init(Amr*)} in \texttt{eib.h}
  \item \texttt{read\_geom()} in \texttt{eib\_io.h}
  \item (later, in AMR init/regrid hooks) \texttt{computeMarkers(lev)}
  \item \texttt{initialiseGPs(lev)}
\end{enumerate}

\subsection{Stage A: \texttt{init} --- Level Metrics}

\texttt{init} performs solver-level setup:
\begin{itemize}[nosep]
  \item Stores \texttt{amr\_p} and refinement ratios.
  \item Allocates per-level containers (\texttt{bmf\_a}, \texttt{faces\_per\_level}).
  \item Computes per-level cell sizes $\Delta x$, cell diagonal $h$, and image-point distances:
  \[
    d^{\mathrm{GP}}_\ell = \alpha\,h_\ell, \qquad
    d^{\mathrm{surf}}_\ell = \alpha_{\mathrm{surf}}\,h_\ell.
  \]
  \item Calls \texttt{read\_geom()}.
\end{itemize}

\subsection{Stage B: \texttt{read\_geom} --- Geometry Load + BVH Build}

Geometry filenames are read from \texttt{ib.filename}.
For each geometry index $i$:

\begin{enumerate}[nosep]
  \item \textbf{Read geometry:}
  \begin{itemize}[nosep]
    \item 2-D: \texttt{read\_polygon\_2d(files\_a[i], geom\_a[i], min\_dx)} ---
          subdivides edges so all segment lengths $\leq \texttt{min\_dx}/2$.
    \item 3-D: \texttt{read\_mesh(files\_a[i], geom\_a[i])} --- reads STL/OFF/VTK.
  \end{itemize}
  \item \textbf{Build BVH:} \texttt{bvh\_a[i].build(geom\_a[i])} --- Morton-sorted median split.
  \item \textbf{Construct inside/outside tester:}
        \texttt{inside\_t(geom\_a[i], bvh\_a[i])} for both 2-D and 3-D.
  \item \textbf{Store bbox:} \texttt{bbox\_a[i] = geom\_a[i].bbox()}.
  \item \textbf{Append flattened per-element cache:}
        \texttt{SurfElem\_a}, \texttt{LocalFrame\_a}, \texttt{geom\_offsets[i]}.
\end{enumerate}

After all geometries: validate cache sizes, run geometry consistency checks
(intersection/containment), optionally flip normals if \texttt{interior\_is\_solid = false}.

\subsection{Stage C: \texttt{computeMarkers} --- Solid/Ghost Classification}

\texttt{computeMarkers(lev)} fills per-cell marker components (comp\,0: fluid/solid id,
comp\,1: ghost-cell id):

\begin{enumerate}[nosep]
  \item Solid pass: \texttt{InsideTesterView} in \texttt{ParallelFor} with bbox rejection
        (GPU path available via AMReX \texttt{AMREX\_USE\_GPU}).
  \item Ghost pass: flags interface-adjacent solids.
  \item Prefix-sum compaction extracts \texttt{gp\_ijk} list.
\end{enumerate}

\subsection{Stage D: \texttt{initialiseGPs} --- Per-Ghost Geometry Metadata}

For each ghost point:

\begin{enumerate}[nosep]
  \item Determine owning geometry index.
  \item BVH closest-point query $\to$ IB foot point and primitive id.
  \item Map primitive id to global face/edge index via \texttt{geom\_offsets}.
  \item Compute distance-to-boundary and local frame.
  \item Place image points and build interpolation stencils/weights.
  \item Store into \texttt{gpData\_t} SoA arrays.
\end{enumerate}

\subsection{Data State after Initialization}

\renewcommand{\arraystretch}{1.2}
\begin{longtable}{@{}L{4.5cm}L{9cm}@{}}
\toprule
\textbf{Object} & \textbf{Description} \\
\midrule
\endfirsthead
\toprule
\textbf{Object} & \textbf{Description} \\
\midrule
\endhead
\texttt{geom\_a}        & Geometry containers (\texttt{TriMesh} or \texttt{Polygon2D}) \\
\texttt{bvh\_a}         & Per-geometry BVH trees \\
\texttt{inout\_fa}      & Inside/outside testers (\texttt{InsideTester*}) \\
\texttt{bbox\_a}        & Coarse geometry bounding boxes \\
\texttt{SurfElem\_a}    & Flattened surface element descriptors \\
\texttt{LocalFrame\_a}  & Flattened normal/tangent basis per element \\
\texttt{geom\_offsets}  & Start index per geometry in flattened arrays \\
\texttt{IBMultiFab} (per level) & Marker field + per-FAB \texttt{gpData\_t} payload \\
\bottomrule
\end{longtable}

% ============================================================
\section{Notes on 2-D vs.\ 3-D}
% ============================================================

\begin{itemize}[nosep]
  \item In 3-D, \texttt{InsideTesterView} uses BVH-guided ray/closest-point logic on
        triangle faces.
  \item In 2-D, \texttt{InsideTesterView} uses BVH-accelerated polygon edge culling for
        both boundary checks and ray-cast inside/outside tests.
        This optimization was added to handle subdivided polygons
        (e.g., 40 $\to$ 736 vertices after \texttt{read\_polygon\_2d}), which caused
        GPU hangs with the previous $O(N)$ linear scan.
  \item Both dimensions share the same dataflow in \texttt{eib\_t}, with
        geometry-specific logic selected by \texttt{AMREX\_SPACEDIM} at compile time.
\end{itemize}

% ============================================================
\section{Primary Runtime Path (after init)}
% ============================================================

\begin{enumerate}[nosep]
  \item \texttt{computeGPs()} --- reconstruct ghost-point flow state from interpolation
        stencils in \texttt{gpData\_t}.
  \item \texttt{computeSURFs()} --- reconstruct surface quantities (velocity, temperature,
        fluxes) at IB surface elements.
  \item \texttt{plotSURF()} --- integrate over \texttt{SurfElem\_a} using precomputed
        \texttt{LocalFrame\_a} and write surface output (drag, heat flux, etc.).
\end{enumerate}

The polygon subdivision in Stage~B ensures sufficient surface point density for accurate
drag and heat flux integration.

\end{document}
